%------------------------------------------------------------------------------%

\usepackage{integracao_e_series}

\title{Teorema Fundamental do Cálculo -- Parte 1}

%------------------------------------------------------------------------------%
\begin{document}

\addtitle

%------------------------------------------------------------------------------%
\section{Teorema Fundamental do Cálculo -- Parte 1}

%------------------------------------------------------------------------------%
\begin{frame}{Teorema Fundamental do Cálculo -- Parte 1}
  Se $f$ for {\color{blue} contínua} em $[a, b]$,
  \pause
  então a função $F$ definida por
  \[
    F(\memph{x})
    \;=\; \int_a^{\memph{x}} f(t)\,dt
    \hspace{15mm} a \leq x \leq b
  \]
  \pause
  é contínua em $[a, b]$,
  \pause
  derivável em $(a, b)$
  \pause
  e
  \[
    \frac{dF}{dx}(x) \;=\; f(x)
  \]
\end{frame}

%------------------------------------------------------------------------------%
\begin{frame}{Em outras palavras}
  Se
  \[
    F(\memph{x})
    \;=\; \int_a^{\memph{x}} f(t)\,dt
  \]
  \pause
  então
  \[
    \frac{d}{d\memph{x}}
    \left[\; \int_a^{\memph{x}} f(t)\,dt \;\right]
    = f(\memph{x})
  \]
\end{frame}

%------------------------------------------------------------------------------%
\section{Demonstração}

%------------------------------------------------------------------------------%
\begin{frame}{Demonstração}
  Queremos mostrar que, se
  \[
    F(x) = \int_a^x f(t)\,dt
  \]
  então
  \[
    \frac{dF}{dx} = f(x)
  \]
\end{frame}

%------------------------------------------------------------------------------%
\begin{frame}{Demonstração}
\onslide<+->{Sejam $x$ e $x+h$ em $(a, b)$}
\begin{align*}
  \onslide<+->{F(x+h) - F(x)
    & = {\color<3>{blue}\int_a^{x+h} f(t)\,dt} - \int_a^x f(t)\,dt                     \\[3mm]}
  \onslide<+->{
    & = {\color<3>{blue}\int_a^x f(t)\,dt + \int_x^{x+h} f(t)\,dt} - \int_a^x f(t)\,dt \\[3mm]}
  \onslide<+->{
    & = \int_x^{x+h} f(t)\,dt}
\end{align*}
\end{frame}

%------------------------------------------------------------------------------%
\begin{frame}{Demonstração}
  Dividindo por $h\neq 0$
  \[
    \frac{F(x+h) - F(x)}{h} = \frac{1}{h} \int_x^{x+h} f(t)\,dt
  \]
  \pause
  No limite $h\to 0$
  \[
    \frac{dF}{dx}
    \pause
    = \lim_{h\to0} \frac{F(x+h) - F(x)}{h}
    \pause
    = \lim_{h\to0} \frac{1}{h} \int_x^{x+h} f(t)\,dt
  \]
  \pause
  Falta mostrar que
  \[
    \lim_{h\to 0} \frac{1}{h} \int_x^{x+h} f(t)\,dt \;\stackrel{?}{=}\; f(x)
  \]
\end{frame}

%------------------------------------------------------------------------------%
\begin{frame}{Demonstração}
  Pelo Teorema Valor Médio,
  \pause
  existe \(c\in[x, x+h]\) tal que
  \[
    \int_x^{x+h} f(t)\,dt = f(c)h
  \]
  \pause
  portanto
  \[
    \frac{1}{h} \int_x^{x+h} f(t)\,dt = f(c)
  \]
  \pause
  Quando $h\to0$,
  \pause
  necessariamente $c\to x$,
  \pause
  assim
  \[
    \frac{dF}{dx}
    \;=\; \lim_{h\to 0}\; \frac{1}{h} \int_x^{x+h} f(t)\,dt
    \;=\; f(x)
  \]
\end{frame}

%------------------------------------------------------------------------------%
\section{Exemplos}

%------------------------------------------------------------------------------%
\addtocounter{exemplo}{1}
\begin{frame}{Exemplo \theexemplo}
  Calcule a derivada da função \(g(x)=\int_0^x \sqrt{1+t^2} \,dt\)
\end{frame}

\begin{frame}{Exemplo \theexemplo}
  Sendo
  \[
    f(t)= \sqrt{1+t^2}
  \]
  contínua em todo o conjunto dos números reais

  \pause
  Pela parte 1 do Teorema Fundamental do Cálculo
  \[
    g'(x)=\sqrt{1+x^2}
  \]
\end{frame}

%------------------------------------------------------------------------------%
\section{Lista Mínima}

%------------------------------------------------------------------------------%
\begin{frame}{Lista Mínima}
  Estudar a Seção 2.5 da Apostila

  \vfill
  Atenção: A prova é baseada no livro, não nas apresentações
\end{frame}

%------------------------------------------------------------------------------%
\end{document}
%------------------------------------------------------------------------------%
